%% Moment d'une force et bras de levier.
%% (a) force perpendiculaire : d = distance PA.
%% (b) force oblique : d se mesure perpendiculairement au SUPPORT de la force.
\documentclass[tikz,border=4pt]{standalone}
\usepackage{liaisons}
\begin{document}
\begin{tikzpicture}[x=1cm,y=1cm,
  force/.style={-{Stealth[length=6pt]}, line width=1.2pt, solideA},
  cote/.style={{Stealth[length=4pt]}-{Stealth[length=4pt]}, line width=0.7pt, solideD},
  support/.style={line width=0.5pt, dash pattern=on 3pt off 2pt, black!55},
  bras/.style={line width=0.6pt, dash pattern=on 3pt off 2pt, black!55}]

%% ============================ (a) ====================================
\begin{scope}[shift={(0,0)}]
  \coordinate (P) at (0,0);   \coordinate (A) at (2.2,0);
  \draw[bras] (P) -- (A);
  \fill (P) circle (2pt);  \fill (A) circle (1.7pt);
  \node[font=\small, anchor=north east, inner sep=2pt] at (P) {$P$};
  \node[font=\small, anchor=south east, inner sep=3pt] at (A) {$A$};
  \draw[force] (A) -- ++(0,1.25) node[anchor=west, inner sep=2pt, font=\small, text=solideA] {$\vec{F}$};
  \draw[line width=0.6pt, black!60] (2.05,0) -- (2.05,0.15) -- (2.2,0.15);
  \draw[cote] (0,-0.55) -- (2.2,-0.55);
  \node[font=\small, text=solideD, anchor=north, inner sep=2pt] at (1.1,-0.58) {$d$};
  \node[font=\small, anchor=north, align=center] at (1.1,-1.62)
       {(a) $\vec{F} \perp PA$\\ $M = F \times d$};
\end{scope}

%% ---------------------------- filet ---------------------------------
\draw[line width=0.4pt, black!35] (3.15,-1.9) -- (3.15,1.5);

%% ============================ (b) ====================================
\begin{scope}[shift={(4.1,0)}]
  \coordinate (P) at (0,0);   \coordinate (A) at (2.2,0);
  % support de F (droite d'action), prolongé de part et d'autre
  \draw[support] ($(A)+(220:1.9)$) -- ($(A)+(40:1.6)$);
  \draw[bras] (P) -- (A);
  \fill (P) circle (2pt);  \fill (A) circle (1.7pt);
  \node[font=\small, anchor=north east, inner sep=2pt] at (P) {$P$};
  \node[font=\small, anchor=north west, inner sep=3pt] at (A) {$A$};
  \draw[force] (A) -- ++(40:1.35) node[anchor=south west, inner sep=1pt, font=\small, text=solideA] {$\vec{F}$};
  % pied de la perpendiculaire abaissée de P sur le support
  \coordinate (H) at (0.909,-1.083);
  \draw[line width=1.1pt, solideD] (P) -- (H);
  \node[font=\small, text=solideD, anchor=west, inner sep=3pt] at (0.5,-0.55) {$d$};
  % marque d'angle droit en H
  \draw[line width=0.6pt, black!60] ($(H)+(40:0.16)$) -- ($(H)+(40:0.16)+(130:0.16)$) -- ($(H)+(130:0.16)$);
  \node[font=\small, anchor=north, align=center] at (1.1,-1.62)
       {(b) $d$ se mesure $\perp$\\ au support de $\vec{F}$};
\end{scope}
\end{tikzpicture}
\end{document}
